\documentclass[conference]{IEEEtran}

% Pacotes essenciais
\usepackage{cite}
\usepackage{graphicx}
\usepackage{amsmath}
\usepackage{float}    % Permite usar o [H] para fixar tabelas
\usepackage{booktabs} % Melhora a estética das tabelas (linhas horizontais bonitas)


% --- Bloco de Preâmbulo (Codificação e Idioma para pdfTeX) ---
% Pacotes para compatibilidade com pdfLaTeX (compilador padrão)
\usepackage[utf8]{inputenc} % Permite caracteres acentuados
\usepackage[T1]{fontenc}    % Codificação de fonte correta

% Configura o idioma principal (português do Brasil)
\usepackage[brazilian]{babel}

% Garante que as listas (itemize) funcionem corretamente
\usepackage{enumitem}
\setlist[itemize]{label=-}



% --- Fim do Bloco ---

% Hyperref é recomendado e deve ser um dos últimos pacotes a serem carregados
\usepackage{hyperref} 

% --- Dados do Artigo ---
\title{ANÁLISE DE VOTAÇÃO PARLAMENTAR COM INTELIGÊNCIA ARTIFICIAL: ADERÊNCIA IDEOLÓGICA}

\author{
    \IEEEauthorblockN{Sarpa, Mauricio}
    \IEEEauthorblockA{\textit{Engenharia da Computação} \\
    \textit{Fundação Hermínio Ometto}\\
    Araras, Brasil\\
    sarpa@alunos.fho.edu.br}
    \and
    \IEEEauthorblockN{Ramos Novais, Pablo}
    \IEEEauthorblockA{\textit{Engenharia da Computação} \\
    \textit{Fundação Hermínio Ometto}\\
    Araras, Brasil\\
    pablo@alunos.fho.edu.br}
}
% -------------------------

\begin{document}

\maketitle

\begin{abstract}
A análise da votação nominal no legislativo é crucial para a transparência. Contudo, a coerência de voto é complexa devido à \textbf{politização} dos temas. Métodos estatísticos tradicionais têm limitações em capturar a natureza multifatorial dessas decisões.Propomos o desenvolvimento de um sistema de Inteligência Artificial para comparar o voto real de deputados e senadores com um voto esperado, inferido por sua classificação política. A principal contribuição é um algoritmo que quantifica a \textbf{aderência ideológica e partidária}, oferecendo uma métrica contínua e objetiva sobre a coerência política no legislativo brasileiro.
\end{abstract}

\begin{IEEEkeywords}
Inteligência Artificial, Ciência Política, Análise de Voto, Politização, Machine Learning.
\end{IEEEkeywords}

\section{Introdução}
\label{sec:introduction}

A análise da \textbf{votação nominal} em casas legislativas é fundamental para a \textbf{transparência política} e a mensuração da ideologia parlamentar. Contudo, a efetiva coerência de voto é frequentemente desafiada pela \textbf{politização} dos temas, que eleva o debate da esfera técnica para a \textbf{esfera ideológica e partidária}, influenciando diretamente a disciplina de voto e o alinhamento esperado \cite{ribeiro2023}.

Os métodos estatísticos tradicionais, como o uso de Componentes Principais \cite{leite2016}, têm limitações em capturar essa complexidade e a natureza multifatorial das decisões legislativas, especialmente com o grande volume de \textbf{dados abertos} disponíveis \cite{lima2023}. Motivado por essa lacuna, este projeto propõe o uso de \textbf{Inteligência Artificial (IA)} e \textit{Machine Learning} para desenvolver ferramentas analíticas mais robustas.

Nosso objetivo central é criar um sistema de IA que \textbf{compare o voto real} dos deputados e senadores com um \textbf{voto esperado}, inferido a partir da sua classificação política e ideológica, mensurando a influência da politização no comportamento legislativo. A principal \textbf{contribuição} deste trabalho reside na criação de um \textbf{algoritmo de classificação} que oferece uma métrica objetiva e contínua para quantificar a \textbf{aderência ideológica/partidária} dos representantes, fornecendo uma nova lente para entender a dinâmica de conflito e a coerência política no legislativo brasileiro.

\section{Trabalhos Relacionados}
\label{sec:related_work}

A análise quantitativa de votações nominais constitui o alicerce para a mensuração da ideologia e da disciplina partidária no legislativo. Historicamente, o \textit{state of the art} nesse campo se concentra em métodos de redução de dimensionalidade para mapear as posições ideológicas dos parlamentares em um espaço de baixa dimensão.

\subsection{Modelos Clássicos e Votações Nominais}
Uma das abordagens mais consolidadas na literatura brasileira é o uso da Análise de Componentes Principais (ACP), ou PCA. Leite e Trento \cite{leite2016} descrevem e validam o PCA para a análise de votações nominais no Brasil, um contexto onde as abstenções e as análises agregadas por partido são cruciais. Os autores demonstram a eficácia do PCA em relação ao modelo W-NOMINATE, estabelecendo o tratamento dos votos (`+1` para sim, `-1` para não, `0` para neutro) como um padrão robusto para a construção de matrizes de votação, o mesmo esquema adotado neste trabalho.

Embora o PCA seja eficaz para posicionar os parlamentares em um espectro ideológico, sua limitação reside na capacidade de capturar a dinâmica da coesão e da indisciplina de voto, especialmente em cenários de alta \textbf{politização} \cite{ribeiro2023}. Buscando aprimorar a identificação desses padrões, Kreuch \cite{kreuch2024} explora o uso de outras técnicas de redução de dimensionalidade (como MDS, t-SNE e UMAP) combinadas com algoritmos de agrupamento mais avançados (como HDBSCAN e Propagação por Afinidade) para identificar a formação de coalizões e a evolução dos agrupamentos ideológicos ao longo do tempo. Esse trabalho reforça a relevância das técnicas de \textit{clustering} para a análise de coesão, premissa fundamental para a metodologia proposta.

\subsection{Análise de Conteúdo e a Lacuna da Aderência}
O debate sobre a verdadeira natureza da ideologia parlamentar levou a estudos que buscam mensurá-la a partir de fontes alternativas. Lima \cite{lima2023} utiliza o Processamento de Linguagem Natural (PLN) e \textit{Machine Learning} para analisar as relações entre partidos e a coesão ideológica baseada em \textbf{discursos} no Congresso Nacional. A motivação para analisar discursos reside na premissa de que a linguagem utilizada pode revelar diferenças intrapartidárias e ideológicas que são mascaradas pela disciplina partidária nos votos nominais.

Neste contexto, o presente trabalho se insere na lacuna identificada pela literatura. Enquanto os modelos existentes focam em \textit{descrever} o voto real (PCA) ou em inferir a ideologia a partir de \textit{discursos} (PLN), inexiste uma métrica contínua e objetiva que \textbf{quantifique a aderência} de um parlamentar ao seu \textbf{voto esperado} ou ideal. Propomos, portanto, um algoritmo que utiliza o K-means não apenas para mapear, mas para \textbf{modelar o voto ideal} de cada agrupamento, permitindo a comparação direta com o voto real do indivíduo. Essa abordagem inova ao fornecer um indicador direto e contínuo da coerência política e da influência da \textbf{politização} na tomada de decisão.

\section{Metodologia Proposta}
\label{sec:metodologia}
A metodologia deste trabalho é estruturada em um pipeline de ciência de dados, compreendendo a aquisição e tratamento dos dados, a aplicação de um modelo de aprendizado não supervisionado para identificação de padrões e, por fim, a validação rigorosa dos resultados obtidos.

\subsection{Coleta e Estruturação dos Dados}
A fonte primária de dados é a API de Dados Abertos do Senado Federal, que disponibiliza o histórico de votações em formato JSON. O processo de coleta é realizado de forma programática, seguindo uma sequência de requisições para obter os votos individuais de cada senador em cada votação. Primeiramente, os dados brutos são transformados em uma \textbf{matriz senador-votação}, na qual cada linha representa um senador e cada coluna, uma votação específica.

Os valores categóricos dos votos ('Sim', 'Não', 'Abstenção', etc.) são convertidos para um formato numérico, seguindo um esquema de codificação consolidado na análise legislativa brasileira \cite{leite2016}:
\begin{itemize}
    \item \textbf{Sim:} +1
    \item \textbf{Não:} -1
    \item \textbf{Abstenção, Obstrução, Ausência ou Presidente:} 0
\end{itemize}

% --- INSERÇÃO: Justificativa da Não-Normalização ---
Diferentemente de abordagens tradicionais que aplicam normalização (como Z-score) aos dados, optou-se por utilizar os valores brutos da codificação proposta. Esta decisão fundamenta-se na necessidade de preservar o valor zero como um ponto fixo de neutralidade (abstenção ou ausência), garantindo que a interpretação geométrica da distância em relação à origem permaneça semânticamente coerente com a posição de não-voto.

Um passo crítico no pré-processamento é a filtragem das votações. Seguindo a recomendação de Ribeiro (2023), apenas votações \textbf{substantivas} (de mérito) são mantidas, enquanto votações puramente \textbf{procedimentais} são descartadas. Essa distinção é crucial para evitar que o modelo agrupe os parlamentares com base em táticas regimentares (ruído) em vez de seus posicionamentos ideológicos (sinal) \cite{ribeiro2023}.

% --- INSERÇÃO: Filtros de Unanimidade e Assiduidade ---
Adicionalmente, foram aplicados dois critérios de filtragem para garantir a qualidade estatística da matriz de dados. Primeiro, foram excluídas as votações unânimes, uma vez que a ausência de variância nesses registros não contribui para a discriminação de agrupamentos (clusters). Segundo, estabeleceu-se um critério de assiduidade mínima: parlamentares com participação inferior a 50\% nas votações do período analisado foram removidos da amostra. Esta medida visa mitigar a esparsidade excessiva dos vetores de votação, que poderia introduzir ruído no cálculo das distâncias euclidianas.

\subsection{Modelo de Agrupamento: K-Means}
Para identificar os agrupamentos latentes de senadores, foi selecionado o algoritmo de clusterização \textbf{K-means}, um método de aprendizado não supervisionado. O K-means particiona o conjunto de dados em $K$ clusters distintos, buscando minimizar a inércia, também conhecida como Soma dos Quadrados Dentro do Cluster (WCSS - \textit{Within-Cluster Sum of Squares}). A inércia é definida como:
$$\sum_{i=0}^{n} \min_{\mu_j \in C} (||x_i - \mu_j||^2)$$
onde $x_i$ é o vetor de votos de um senador e $\mu_j$ é o centróide (ponto médio) do cluster $C_j$.

O algoritmo opera iterativamente em duas etapas:
\begin{enumerate}
    \item \textbf{Atribuição:} Cada senador (ponto de dado) é atribuído ao cluster cujo centróide é o mais próximo, com base na distância euclidiana.
    \item \textbf{Atualização:} O centróide de cada cluster é recalculado como a média dos vetores de todos os senadores atribuídos a ele.
\end{enumerate}
Essas etapas são repetidas até a convergência, quando as atribuições de cluster não se alteram mais. Para a inicialização dos centróides, utiliza-se o método \textit{K-means++}, que melhora a qualidade dos resultados e acelera a convergência.

\subsection{Validação e Otimização do Modelo}
A escolha do número ótimo de clusters ($K$) é um parâmetro fundamental e não arbitrário.
% --- SUBSTITUIÇÃO: Validação por Ensemble ---
A determinação do número ótimo de clusters ($K$) não se baseou em uma única métrica. Adotou-se uma abordagem de validação cruzada por múltiplos índices para mitigar vieses específicos de cada algoritmo. Além do Método do Cotovelo e da Análise de Silhueta, foram calculados os índices Davies-Bouldin (onde menores valores indicam melhor separação) e Calinski-Harabasz (onde maiores valores indicam clusters mais densos).

O valor final de $K$ foi definido através de um sistema de consenso (votação majoritária) entre as métricas sugeridas por cada índice, selecionando o número de agrupamentos que apresentou maior estabilidade entre os diferentes métodos de avaliação.

\subsection{Algoritmo de Classificação para Aderência Ideológica}
A principal contribuição deste trabalho reside no desenvolvimento de um algoritmo que transforma a saída do modelo K-means (o agrupamento ideológico) em uma métrica de \textbf{Aderência Ideológica} ($\text{AI}$) contínua e interpretável.

\subsubsection{Definição do Voto Esperado}
Após a clusterização da matriz senador-votação pelo K-means, cada um dos $K$ clusters ($C_j$) representa uma \textbf{posição ideológica latente}. O centróide do cluster ($\mu_j$) corresponde ao vetor de votos médios do grupo em todas as votações. Para definir o \textbf{Voto Esperado} ($\hat{v}_{j}$) desse grupo, o centróide é binarizado (ou trinarizado) por meio de um limiar ($\tau$).

O \textbf{Voto Esperado do Cluster} ($\hat{v}_{j}^{(i)}$) para uma votação $i$ é dado por:
$$
\hat{v}_{j}^{(i)} = \begin{cases} 
+1 & \text{se } \mu_j^{(i)} > \tau \\ 
-1 & \text{se } \mu_j^{(i)} < -\tau \\ 
0 & \text{se } -\tau \le \mu_j^{(i)} \le \tau 
\end{cases}
$$
% --- ATUALIZAÇÃO: Tau = 0.1 ---
onde $\tau$ é um limiar de neutralidade definido empiricamente neste estudo como $\tau=0.1$. Este limiar estabelece uma zona de neutralidade mais robusta, assegurando que apenas tendências majoritárias claras no cluster sejam classificadas como Voto Esperado (+1 ou -1), enquanto médias próximas a zero (indicativas de dispersão interna no grupo) permanecem classificadas como neutras.

\subsubsection{Cálculo da Aderência Ideológica ($\text{AI}$)}
A métrica $\text{AI}$ quantifica a coerência do senador $i$ em relação à posição ideológica ideal de seu grupo $C_{\text{ideal}}$. Essa aderência é calculada utilizando a \textbf{Similaridade de Cosseno} entre o vetor de voto real do senador ($x_i$) e o vetor de voto esperado de seu cluster ($\hat{v}_{\text{ideal}}$), sendo o cluster $C_{\text{ideal}}$ aquele ao qual o senador foi atribuído pelo K-means.

A Similaridade de Cosseno é definida como:
$$
\text{AI}_i = \text{Similaridade}(\mathbf{x}_i, \mathbf{\hat{v}}_{\text{ideal}}) = \frac{\mathbf{x}_i \cdot \mathbf{\hat{v}}_{\text{ideal}}}{\|\mathbf{x}_i\| \|\mathbf{\hat{v}}_{\text{ideal}}\|}
$$
O resultado é um valor contínuo que varia de -1 (totalmente oposto ao voto esperado) a +1 (aderência perfeita), oferecendo um índice robusto de quão bem o parlamentar se alinha ao seu agrupamento ideológico.

\section{Resultados e Discussão}

A aplicação do pipeline de orquestração permitiu processar os dados de votações nominais do Senado Federal referente ao ano de 2023. A seguir, detalham-se as etapas de redução de dimensionalidade, a determinação do número ótimo de clusters e a análise aprofundada da aderência ideológica.

\subsection{Pré-processamento e Redução de Dimensionalidade}

A coleta inicial retornou 141 votações. Após a filtragem metodológica — removendo votações de homenagens, pesar e aquelas com unanimidade (baixa variância) —, restaram 35 votações significativas para a análise. O dataset final contou com 81 senadores ativos (participação superior a 50\%).

A análise descritiva revelou uma esparsidade de 17,60\%, indicando um preenchimento consistente da matriz de votos. A aplicação da Análise de Componentes Principais (PCA) reduziu as 35 dimensões originais (votações) para 19 componentes principais, preservando 95\% da variância explicada dos dados.

\subsection{Determinação do Número de Clusters (K)}

Para definir o número ideal de agrupamentos, o algoritmo K-Means foi executado com $K$ variando de 2 a 10. As métricas de validação são detalhadas na Tabela \ref{tab:metrics} e visualizadas na Figura \ref{fig:validacao}.

\begin{table}[H]
\caption{Métricas de Validação para Diferentes Valores de K}
\begin{center}
\begin{tabular}{|c|c|c|c|}
\hline
\textbf{K} & \textbf{Inércia} & \textbf{Silhueta} & \textbf{Davies-Bouldin} \\
\hline
2 & 1671.65 & \textbf{0.3205} & \textbf{1.2530} \\
3 & 1375.97 & 0.2439 & 1.4711 \\
4 & 1253.96 & 0.1948 & 1.7596 \\
5 & 1190.74 & 0.1977 & 1.7831 \\
\hline
\end{tabular}
\label{tab:metrics}
\end{center}
\end{table}

\begin{figure}[H]
\centerline{\includegraphics[width=0.48\textwidth]{validacao_k_otimo_completa.png}}
\caption{Comportamento das métricas de validação. Nota-se que K=3 mantém estabilidade antes da queda acentuada de desempenho em K=4.}
\label{fig:validacao}
\end{figure}

A análise da Figura \ref{fig:validacao} revela que, embora $K=2$ apresente os melhores indicadores absolutos (sugerindo polarização binária), a transição para $K=3$ não representa uma perda estrutural crítica. O índice de Silhueta em $K=3$ (0.2439) permanece significativamente superior ao patamar observado em $K=4$ (0.1948), onde ocorre uma degradação acentuada da coesão dos grupos.

Portanto, $K=3$ configura-se como um ponto de inflexão: ele oferece uma redução substancial na Inércia (de 1671 para 1375) em comparação a $K=2$, sem sofrer a pulverização excessiva observada em $K \ge 4$. Essa característica matemática valida a escolha manual de três clusters, permitindo capturar a complexidade do espectro político (Governo, Oposição e Centro) que seria ignorada em uma abordagem puramente binária.




\subsection{Caracterização dos Agrupamentos}

Com $K=3$, os perfis foram definidos como:
\begin{itemize}
    \item \textbf{Cluster 0 (Governo):} 25 membros, alta coesão (Aderência média: 0.471). Predominância de PT, PSD e MDB.
    \item \textbf{Cluster 1 (Oposição):} 27 membros, média coesão (Aderência média: 0.419). Predominância de PL, PP e Republicanos.
    \item \textbf{Cluster 2 (Centro/Independente):} 29 membros, baixa coesão (Aderência média: 0.258). Grupo heterogêneo com membros de PSD, União e Podemos.
\end{itemize}

\subsection{Análise de Aderência Ideológica (AI)}

A métrica de Aderência Ideológica (AI) permitiu classificar os senadores quanto à fidelidade ao comportamento médio de seu grupo.

A Tabela \ref{tab:top_aligned} apresenta os 15 parlamentares com maior aderência. Nota-se que senadores do PP e PL dominam o topo da lista de fidelidade no Cluster 1, enquanto o PT e aliados demonstram forte alinhamento no Cluster 0. O senador Luis Carlos Heinze (PP) obteve a maior aderência absoluta (0.7588).

\begin{table}[H]
\caption{Top 15 Parlamentares com Maior Aderência}
\begin{center}
\begin{tabular}{|l|c|c|c|}
\hline
\textbf{Parlamentar} & \textbf{Partido} & \textbf{Cluster} & \textbf{Aderência (AI)} \\
\hline
Luis Carlos Heinze & PP & 1 & 0.7588 \\
Eliziane Gama & PSD & 0 & 0.6857 \\
Wilder Morais & PL & 1 & 0.6583 \\
Tereza Cristina & PP & 1 & 0.6564 \\
Beto Faro & PT & 0 & 0.6560 \\
Randolfe Rodrigues & REDE & 0 & 0.6560 \\
Fabiano Contarato & PT & 0 & 0.6560 \\
Paulo Paim & PT & 0 & 0.6560 \\
Jussara Lima & PSD & 0 & 0.6534 \\
Rogério Carvalho & PT & 0 & 0.6527 \\
Augusta Brito & PT & 0 & 0.6477 \\
Rogerio Marinho & PL & 1 & 0.6358 \\
Romário & PL & 1 & 0.6209 \\
Professora Dorinha Seabra & UNIÃO & 2 & 0.5852 \\
Marcelo Castro & MDB & 0 & 0.5735 \\
\hline
\end{tabular}
\label{tab:top_aligned}
\end{center}
\end{table}

Por outro lado, a Tabela \ref{tab:least_aligned} destaca os parlamentares com menor aderência. É notável a presença massiva de membros do Cluster 2 (Centro) nesta lista. Senadores como Oriovisto Guimarães e Laércio Oliveira apresentaram valores de aderência negativos, indicando que seus votos frequentemente divergem do centroide do seu próprio grupo, validando a hipótese de que o "Centro" não opera como um bloco monolítico, mas sim como uma zona de convergência flutuante.

\begin{table}[H]
\caption{Top 15 Parlamentares com Menor Aderência}
\begin{center}
\begin{tabular}{|l|c|c|c|}
\hline
\textbf{Parlamentar} & \textbf{Partido} & \textbf{Cluster} & \textbf{Aderência (AI)} \\
\hline
Oriovisto Guimarães & PODEMOS & 2 & -0.1426 \\
Laércio Oliveira & PP & 2 & -0.1348 \\
Vanderlan Cardoso & PSD & 2 & -0.1241 \\
Alan Rick & UNIÃO & 2 & -0.0243 \\
Soraya Thronicke & UNIÃO & 2 & 0.0045 \\
Zequinha Marinho & PL & 2 & 0.0252 \\
Rodrigo Pacheco & PSD & 2 & 0.0745 \\
Carlos Viana & PODEMOS & 2 & 0.0921 \\
Styvenson Valentim & PODEMOS & 2 & 0.0938 \\
Alessandro Vieira & PSDB & 2 & 0.1268 \\
Marcio Bittar & UNIÃO & 1 & 0.1372 \\
Marcos Rogério & PL & 1 & 0.1794 \\
Marcos do Val & PODEMOS & 1 & 0.1836 \\
Mara Gabrilli & PSD & 2 & 0.1862 \\
Sergio Moro & UNIÃO & 1 & 0.2130 \\
\hline
\end{tabular}
\label{tab:least_aligned}
\end{center}
\end{table}



\section{Conclusão}
\label{sec:conclusao}

O presente trabalho atingiu seu objetivo central ao desenvolver e aplicar uma metodologia baseada em Inteligência Artificial para mensurar a aderência ideológica no Senado Federal. A utilização do algoritmo K-Means, combinada com métricas de validação e uma análise qualitativa do cenário político, permitiu transcender a análise puramente estatística, oferecendo uma interpretação computacional da dinâmica parlamentar.

Os resultados obtidos ratificam a hipótese de que a polarização política brasileira, embora matematicamente binária (como sugerido pelos índices de Silhueta para $K=2$), exige uma modelagem tricotômica para representar a realidade fática. A decisão estratégica de fixar $K=3$ revelou-se acertada ao isolar um terceiro agrupamento ("Centro") cuja característica definidora não é a coesão, mas a dispersão. Enquanto os clusters de Governo e Oposição apresentaram altas taxas de aderência ideológica (médias de 0.47 e 0.41, respectivamente), o Cluster 2 demonstrou uma aderência significativamente menor (0.25) e diversos valores negativos. Isso valida computacionalmente o conceito de "Centrão" como uma zona de convergência flutuante, e não como um bloco ideológico monolítico.

A principal contribuição técnica deste estudo é a formalização da métrica de \textbf{Aderência Ideológica (AI)}. Ao quantificar a fidelidade do parlamentar em uma escala contínua, o sistema oferece uma ferramenta objetiva para identificar tanto os líderes de núcleo duro (como demonstrado no Cluster 1) quanto os parlamentares pendulares.

\subsection{Limitações e Desafios}
Uma limitação inerente ao estudo foi definir o limiar de neutralidade ($\tau$) e a escolha do número de clusters exigiu uma intervenção humana supervisionada sobre o modelo não supervisionado. Embora matematicamente justificada pela estabilidade da silhueta em $K=3$, essa etapa sublinha que a IA, neste contexto, atua melhor como uma ferramenta de suporte à decisão do analista político do que como um oráculo totalmente autônomo.

\subsection{Trabalhos Futuros}
Para extensões futuras, sugere-se a aplicação desta metodologia em uma série temporal longitudinal, abrangendo múltiplas legislaturas, para visualizar a migração de parlamentares entre clusters (mudança de aderência) ao longo do tempo. Outra via promissora é a integração de técnicas de Processamento de Linguagem Natural (PLN) para correlacionar a aderência ideológica votada (o que o senador faz) com a aderência discursiva (o que o senador fala em tribuna), permitindo identificar dissonâncias entre o discurso e a prática legislativa.

Em suma, a ferramenta desenvolvida demonstra que a aplicação de aprendizado de máquina em dados governamentais abertos é um caminho viável e robusto para aumentar a transparência e a compreensibilidade do complexo jogo democrático.

% --- REFERÊNCIAS ---
\begin{thebibliography}{9}

\bibitem{ribeiro2023}
RIBEIRO, T. D. \textit{Por que um deputado é indisciplinado?: uma análise das votações nominais na câmara dos deputados entre os anos de 2015 e 2022}. 2023. Dissertação (Mestrado em Administração Pública) – Instituto Brasileiro de Ensino, Desenvolvimento e Pesquisa, Brasília, 2023.

\bibitem{leite2016}
LEITE, L.; TRENTO, S. Análise de votações nominais do legislativo brasileiro utilizando componentes principais. \textit{Leviathan: Cadernos de Pesquisa Política}, n. 12, p. 120-163, 2016.

\bibitem{lima2023}
LIMA, W. P. C. \textit{Uma análise de partidos políticos baseada em discursos no congresso nacional brasileiro}. 2023. Dissertação (Mestrado em Ciência da Computação) – Centro Federal de Educação Tecnológica Celso Suckow da Fonseca, Rio de Janeiro, 2023.

\bibitem{rousseeuw1987}
ROUSSEEUW, P. J. Silhouettes: a Graphical Aid to the Interpretation and Validation of Cluster Analysis. \textit{Journal of Computational and Applied Mathematics}, v. 20, p. 53-65, 1987.

\bibitem{kreuch2024}
KREUCH, M. A. \textit{Desenvolvimento de solução para visualização e análise de votações políticas}. 2024. Trabalho de Conclusão de Curso (Graduação em Ciências da Computação) – Universidade Federal de Santa Catarina, Florianópolis, 2024.

\end{thebibliography}

\end{document}